\chapter{一章学会建模与训练}
本章把 PyTorch 建模与训练的高频知识压成一条主线：写模型、选层、选损失、
接优化器、喂数据，最后套标准训练循环。面试时只要抓住这条链路，多数问题都能
回答得清楚：数据进入模型，模型输出 logits，损失计算标量，反向传播产生梯度，
优化器根据梯度更新参数。
\section{nn.Module：自定义模型}
自定义模型通常继承 \code{nn.Module}。在 \code{\_\_init\_\_} 中声明层，在
\code{forward} 中描述张量流动。调用 \code{model(x)} 时会自动执行
\code{forward}，不要手动写 \code{model.forward(x)}。
\begin{lstlisting}
import torch
from torch import nn
class MLP(nn.Module):
    def __init__(self, in_dim, hidden_dim, num_classes):
        super().__init__()
        self.fc1 = nn.Linear(in_dim, hidden_dim)
        self.act = nn.ReLU()
        self.drop = nn.Dropout(p=0.2)
        self.fc2 = nn.Linear(hidden_dim, num_classes)
    def forward(self, x):
        x = self.fc1(x)
        x = self.act(x)
        x = self.drop(x)
        return self.fc2(x)
model = MLP(20, 64, 3)
x = torch.randn(8, 20)
logits = model(x)
print(logits.shape)
\end{lstlisting}
\subsection{parameters、named\_parameters 与 device}
\code{model.parameters()} 返回模型参数迭代器，常直接传给优化器。
\code{model.named\_parameters()} 会同时返回参数名，适合调试、冻结层或给不同层设置
不同学习率。
\begin{lstlisting}
for name, param in model.named_parameters():
    print(name, tuple(param.shape), param.requires_grad)
optimizer = torch.optim.AdamW(model.parameters(), lr=1e-3)
\end{lstlisting}
模型和数据必须在同一设备。标准写法是先确定 device，再把模型、输入和 target
都移动过去。
\begin{lstlisting}
device = torch.device("cuda" if torch.cuda.is_available() else "cpu")
model = MLP(20, 64, 3).to(device)
x = torch.randn(8, 20).to(device)
y = torch.randint(0, 3, (8,), device=device)
logits = model(x)
\end{lstlisting}
\subsection{train 与 eval}
\code{model.train()} 与 \code{model.eval()} 切换模块行为，不负责开关梯度。
它们主要影响 Dropout 和 BatchNorm；关闭梯度应使用 \code{torch.no\_grad()}。
\begin{center}
\begin{tabular}{lll}
\toprule
模式 & Dropout & BatchNorm \\
\midrule
\code{train()} & 随机丢弃神经元 & 使用当前 batch 统计量并更新 running 统计量 \\
\code{eval()} & 关闭随机丢弃 & 使用训练期累计的 running 统计量 \\
\bottomrule
\end{tabular}
\end{center}
因此训练循环中先写 \code{model.train()}；验证和测试时写
\code{model.eval()}，并把推理代码包在 \code{with torch.no\_grad():} 中。
\section{常用层速览}
\subsection{Linear、Conv2d 与激活函数}
\code{nn.Linear(in\_features, out\_features)} 是全连接层，输入最后一维必须等于
\code{in\_features}。\code{nn.Conv2d} 常用于图像，输入形状是 \code{[N, C, H, W]}。
常见激活有 \code{nn.ReLU}、\code{nn.GELU}、\code{nn.Sigmoid} 和 \code{nn.Tanh}。
\begin{lstlisting}
linear = nn.Linear(128, 10)
x = torch.randn(32, 128)
print(linear(x).shape)
conv = nn.Conv2d(3, 16, kernel_size=3, padding=1)
image = torch.randn(4, 3, 32, 32)
print(conv(image).shape)
\end{lstlisting}
\subsection{nn.Sequential}
\code{nn.Sequential} 适合顺序堆叠网络层。若模型有多输入、多分支、条件逻辑或跳连，
应直接在 \code{forward} 中写清楚。
\begin{lstlisting}
cnn = nn.Sequential(
    nn.Conv2d(3, 16, kernel_size=3, padding=1),
    nn.ReLU(),
    nn.MaxPool2d(kernel_size=2),
    nn.Conv2d(16, 32, kernel_size=3, padding=1),
    nn.ReLU(),
    nn.AdaptiveAvgPool2d((1, 1)),
    nn.Flatten(),
    nn.Linear(32, 10),
)
\end{lstlisting}
\section{损失函数}
\subsection{MSELoss}
\code{nn.MSELoss} 常用于回归，模型输出和 target 通常形状一致。
\begin{lstlisting}
criterion = nn.MSELoss()
pred = torch.randn(16, 1)
target = torch.randn(16, 1)
loss = criterion(pred, target)
\end{lstlisting}
\subsection{CrossEntropyLoss}
\code{nn.CrossEntropyLoss} 用于单标签多分类。它内部已经包含 \code{LogSoftmax}
和负对数似然损失，所以输入必须是 logits，不要手动 softmax。target 是类别索引，
形状通常为 \code{[N]}，类型为 \code{torch.long}，不是 one-hot。
\begin{lstlisting}
criterion = nn.CrossEntropyLoss()
logits = torch.randn(8, 5)          # raw scores, not probabilities
target = torch.tensor([0, 1, 4, 3, 2, 1, 0, 4])
loss = criterion(logits, target)
\end{lstlisting}
\subsection{BCEWithLogitsLoss}
\code{nn.BCEWithLogitsLoss} 用于二分类或多标签分类，内部包含 Sigmoid 和二元交叉熵，
比先 \code{sigmoid} 再 \code{BCELoss} 更稳定。target 通常是与 logits 同形状的浮点张量。
\begin{lstlisting}
criterion = nn.BCEWithLogitsLoss()
logits = torch.randn(8, 4)
target = torch.randint(0, 2, (8, 4)).float()
loss = criterion(logits, target)
\end{lstlisting}
\section{优化器与学习率调度}
\code{optim.SGD} 简单稳定，配合 momentum 常用于视觉任务；\code{Adam} 自适应调整
每个参数的步长，收敛初期通常更快；\code{AdamW} 将 weight decay 与 Adam 更新解耦，
是 Transformer 等模型的常用选择。\code{weight\_decay} 用来抑制参数过大，降低过拟合。
\begin{lstlisting}
optimizer = torch.optim.SGD(model.parameters(), lr=0.1, momentum=0.9)
optimizer = torch.optim.Adam(model.parameters(), lr=1e-3)
optimizer = torch.optim.AdamW(model.parameters(), lr=1e-3, weight_decay=0.01)
\end{lstlisting}
\code{StepLR} 按固定间隔衰减学习率；\code{CosineAnnealingLR} 按余弦曲线平滑下降。
warmup 指训练初期从小学习率逐步升到目标学习率，常用于大模型或大 batch 训练。
\begin{lstlisting}
optimizer = torch.optim.AdamW(model.parameters(), lr=1e-3)
scheduler = torch.optim.lr_scheduler.StepLR(optimizer, step_size=10, gamma=0.1)
for epoch in range(30):
    train_one_epoch()
    scheduler.step()
cosine = torch.optim.lr_scheduler.CosineAnnealingLR(optimizer, T_max=30)
\end{lstlisting}
\section{Dataset 与 DataLoader}
自定义 \code{Dataset} 至少实现 \code{\_\_len\_\_} 和 \code{\_\_getitem\_\_}。
\code{DataLoader} 负责 batch 化、打乱、并行加载和拼接样本。
\begin{lstlisting}
from torch.utils.data import Dataset, DataLoader
class TensorDataset2D(Dataset):
    def __init__(self, x, y):
        self.x = x
        self.y = y
    def __len__(self):
        return len(self.x)
    def __getitem__(self, index):
        return self.x[index], self.y[index]
dataset = TensorDataset2D(torch.randn(100, 20), torch.randint(0, 3, (100,)))
loader = DataLoader(dataset, batch_size=32, shuffle=True, num_workers=0)
\end{lstlisting}
\begin{center}
\begin{tabular}{ll}
\toprule
参数 & 作用 \\
\midrule
\code{batch\_size} & 每个 batch 的样本数 \\
\code{shuffle} & 训练集通常为 True，验证集通常为 False \\
\code{num\_workers} & 子进程加载数据的数量 \\
\code{drop\_last} & 是否丢弃最后一个不足 batch 的小批次 \\
\code{collate\_fn} & 自定义样本如何拼成 batch，常用于变长序列 \\
\bottomrule
\end{tabular}
\end{center}
\section{完整训练循环模板}
标准五步是 \code{optimizer.zero\_grad()} $\rightarrow$ forward $\rightarrow$ loss
$\rightarrow$ \code{loss.backward()} $\rightarrow$ \code{optimizer.step()}。
下面代码可直接运行，模拟一个三分类任务。
\begin{lstlisting}
import torch
from torch import nn
from torch.utils.data import DataLoader, Dataset
class ToyDataset(Dataset):
    def __init__(self, n=512, in_dim=20, num_classes=3):
        self.x = torch.randn(n, in_dim)
        w = torch.randn(in_dim, num_classes)
        self.y = (self.x @ w).argmax(dim=1)
    def __len__(self):
        return len(self.x)
    def __getitem__(self, index):
        return self.x[index], self.y[index]
class Classifier(nn.Module):
    def __init__(self):
        super().__init__()
        self.net = nn.Sequential(
            nn.Linear(20, 64),
            nn.ReLU(),
            nn.Dropout(p=0.1),
            nn.Linear(64, 3),
        )
    def forward(self, x):
        return self.net(x)
def train_one_epoch(model, loader, criterion, optimizer, device):
    model.train()
    total_loss, total_correct, total_count = 0.0, 0, 0
    for x, y in loader:
        x, y = x.to(device), y.to(device)
        optimizer.zero_grad()
        logits = model(x)
        loss = criterion(logits, y)
        loss.backward()
        optimizer.step()
        batch_size = x.size(0)
        total_loss += loss.item() * batch_size
        total_correct += (logits.argmax(dim=1) == y).sum().item()
        total_count += batch_size
    return total_loss / total_count, total_correct / total_count
def evaluate(model, loader, criterion, device):
    model.eval()
    total_loss, total_correct, total_count = 0.0, 0, 0
    with torch.no_grad():
        for x, y in loader:
            x, y = x.to(device), y.to(device)
            logits = model(x)
            loss = criterion(logits, y)
            batch_size = x.size(0)
            total_loss += loss.item() * batch_size
            total_correct += (logits.argmax(dim=1) == y).sum().item()
            total_count += batch_size
    return total_loss / total_count, total_correct / total_count
def main():
    device = torch.device("cuda" if torch.cuda.is_available() else "cpu")
    train_loader = DataLoader(ToyDataset(512), batch_size=32, shuffle=True)
    val_loader = DataLoader(ToyDataset(128), batch_size=64, shuffle=False)
    model = Classifier().to(device)
    criterion = nn.CrossEntropyLoss()
    optimizer = torch.optim.AdamW(model.parameters(), lr=1e-3, weight_decay=0.01)
    scheduler = torch.optim.lr_scheduler.StepLR(optimizer, step_size=5, gamma=0.5)
    for epoch in range(10):
        train_loss, train_acc = train_one_epoch(
            model, train_loader, criterion, optimizer, device
        )
        val_loss, val_acc = evaluate(model, val_loader, criterion, device)
        scheduler.step()
        print(
            f"epoch={epoch + 1} train_loss={train_loss:.4f} "
            f"train_acc={train_acc:.4f} val_loss={val_loss:.4f} "
            f"val_acc={val_acc:.4f}"
        )
if __name__ == "__main__":
    main()
\end{lstlisting}
\begin{trap}
三类坑最常见：忘记 \code{optimizer.zero\_grad()} 会让梯度跨 batch 累加；验证时
忘记 \code{model.eval()} 会让 Dropout 和 BatchNorm 仍按训练模式工作；统计 loss
时应累加 \code{loss.item()}，否则可能保留计算图并导致显存增长。
\end{trap}
\section{面试高频问答}
\begin{interview}
使用 \code{CrossEntropyLoss} 前要先 softmax 吗？
\end{interview}
\begin{answer}
不要。\code{nn.CrossEntropyLoss} 内部已经包含 \code{LogSoftmax} 和 NLLLoss，
输入应是 logits。手动 softmax 会降低数值稳定性，也可能影响梯度。target 应是类别索引。
\end{answer}
\begin{interview}
为什么必须 \code{zero\_grad()}？
\end{interview}
\begin{answer}
PyTorch 默认累加梯度，而不是自动覆盖梯度。普通训练中如果不清零，本 batch 梯度会叠加
历史 batch 梯度，导致更新方向和步长错误。只有明确做梯度累积时才故意不每步清零。
\end{answer}
\begin{interview}
\code{model.train()} 与 \code{model.eval()} 有什么区别？
\end{interview}
\begin{answer}
它们切换模块行为，主要影响 Dropout 和 BatchNorm。训练模式启用 Dropout 并更新
BatchNorm running 统计量；推理模式关闭 Dropout，并使用累计统计量。关闭梯度要靠
\code{torch.no\_grad()}。
\end{answer}
\begin{interview}
Adam 和 SGD 如何取舍？
\end{interview}
\begin{answer}
Adam 收敛快、调参省心，适合快速建立 baseline；SGD 加 momentum 简单稳定，在部分视觉
任务中泛化好但更依赖学习率调度。现代项目常优先试 AdamW，再按任务验证 SGD。
\end{answer}
\begin{interview}
\code{DataLoader} 的 \code{num\_workers} 是什么？越大越好吗？
\end{interview}
\begin{answer}
它控制用多少子进程加载数据，可减少 GPU 等数据的时间。但不是越大越好，过大会增加
进程切换、内存占用和 I/O 竞争。调试或 Windows 环境常用 0，训练时逐步测试 2、4、8。
\end{answer}
\section{章末速记}
\begin{quickcard}
\textbf{标准训练循环模板}
\begin{enumerate}
  \item \code{model.to(device)}，每个 batch 的 \code{x} 和 \code{y} 也要到同一 device。
  \item 训练前 \code{model.train()}。
  \item 五步：\code{optimizer.zero\_grad()} $\rightarrow$ \code{logits = model(x)}
  $\rightarrow$ \code{loss = criterion(logits, y)} $\rightarrow$
  \code{loss.backward()} $\rightarrow$ \code{optimizer.step()}。
  \item 验证前 \code{model.eval()}，外层写 \code{with torch.no\_grad():}。
  \item \code{CrossEntropyLoss} 吃 logits，target 是类别索引；统计 loss 用
  \code{loss.item()}。
\end{enumerate}
\end{quickcard}
